% 01_motivation.tex -- Part 1, motivation. Sources: research_monograph/chapters/M01_introduction.tex (scaling, BEOL, In2O3
% at the thickness limit, research problem, approach), M03_literature_materials.tex (W and amorphous-oxide literature),
% M10_w_and_vacancies.tex (C023: the paper's crystalline DFT did not test vacancy passivation), M14_audit.tex (published
% crystalline PBE of In2O3 and W:In2O3), M02_questions.tex and M08_c10_campaign.tex (register, shared model, stages
% S0-S7, S0/S1 seen before the freeze, donor density fixed at the model V1 value), M04_dft_theory.tex (HSE06 checks),
% M12_multiscale.tex (transfer table: W treated as fully ionised, no donor density transferred), M15_conclusions.tex
% (6.3 nm out-of-sample misses), 03_device_data_rev.tex and M07_provenance.tex (published mobilities, pp. 3, 7),
% 11_validation_rev.tex (per-film misses), evidence/FACT_REGISTRY.tsv (FACT-101..104), CLAIM_SOURCE_MAP.tsv (C002, C010,
% C012-C014, C016, C023, C029, M007: 2.0-13.2 nm step -0.579 V, 2.0 -> 6.3 nm -0.018 V).
% Deck rules (AUTHORING_GUIDE section 8): step terminology G4 is defined on the problem slide; at most three footnote
% lines per slide (G6), so works cited for consecutive statements share one footnote; no \sref inside \src (G3).
\providecommand{\oneTop}{\vspace{-3mm}}

\partslide{Part 1}{Motivation: ultrathin oxide channels}

% ---------------------------------------------------------------------------------------------------------------------
\begin{frame}{Upper transistor tiers must be processed below about 400\,\textdegree C}
  \oneTop
  \begin{columns}[T,onlytextwidth]
    \begin{column}{0.60\textwidth}
      \begin{itemize}
        \item Geometric scaling has slowed since the end of Dennard scaling, and thinned channels lose performance:
          Si below about 3\,nm, InGaAs below about 10\,nm\footnote[frame]{\refcite{Theis2017}; \refcite[p.~164]{Si2022}}.
        \item Monolithic 3-D integration stacks back-end-of-line (BEOL) transistor tiers above the Cu interconnect,
          which limits their processing to below about 400\,\textdegree C\sref[p.~1]{Januar2026}.
        \item This excludes crystalline Si; oxide semiconductors are the leading candidates. Their low off-current
          enables capacitor-less DRAM with retention beyond 400\,s\srefs{Datta2019,Belmonte2020}.
      \end{itemize}
    \end{column}
    \begin{column}{0.36\textwidth}
      \centering
      \vspace{1mm}
      \begin{tikzpicture}[x=1mm,y=1mm,line width=0.5pt,font=\scriptsize]
        \draw (0,0) rectangle (36,6);   \node at (18,3) {Si substrate};
        \draw (0,6) rectangle (36,13);  \node at (18,9.5) {Si CMOS logic};
        \draw (0,13) rectangle (36,29);
        \foreach \y in {16.5,21,25.5} {\draw[inkgrey] (3,\y) -- (10,\y); \draw[inkgrey] (26,\y) -- (33,\y);}
        \foreach \x in {6,30} {\draw[inkgrey] (\x,13) -- (\x,29);}
        \node[align=center] at (18,21) {Cu\\interconnect};
        \draw[ntumaroon,line width=0.8pt] (0,29) rectangle (36,37);
        \node[align=center] at (18,33) {oxide-channel\\transistor tier};
        \draw[inkgrey] (0,37) rectangle (36,42); \node[inkgrey] at (18,39.5) {upper interconnect};
        \draw (37,42) -- (38.5,42) -- (38.5,29) -- (37,29);
        \node[anchor=west,align=left] at (39,35.5) {below\\400\,\textdegree C};
      \end{tikzpicture}
      \src{Schematic of a monolithic 3-D stack, not to scale.}
    \end{column}
  \end{columns}
  \takeaway{Oxide semiconductors are the leading channel candidates for the upper tiers.}
\end{frame}

% ---------------------------------------------------------------------------------------------------------------------
\begin{frame}{Indium oxide switches off only as a film a few nanometres thick}
  \oneTop
  \begin{columns}[T,onlytextwidth]
    \begin{column}{0.46\textwidth}
      {\small\textbf{Why oxide and \InO{} channels}\par}
      \begin{itemize}
        \item Amorphous oxides were developed for display transistors; their s-like conduction band survives
          disorder\srefs{Nomura2004,Medvedeva2017}.
        \item \InO{} switches off only when very thin: atomic-layer-deposited 0.7\,nm channels work in enhancement mode,
          with a threshold tuned by thickness and attributed to confinement. Nanometre-thick \InO{} carries drain
          currents above 10\,A/mm\footnote[frame]{\refcite[p.~500]{Si2021}; \refcite[p.~21536]{Lin2022}}.
      \end{itemize}
    \end{column}
    \begin{column}{0.50\textwidth}
      {\small\textbf{What thinning costs}\par}
      \begin{itemize}
        \item Confinement opens the gap and moves the band edges (Part~6) \lab{computed}.
        \item Every atom is near a surface; conventional amorphous oxides become disorder-limited below about
          10\,nm\sref[p.~1]{Januar2026}.
        \item Device spread grows as donors become few: about 500 per \textmu m$^{2}$ in a 2.0\,nm channel at the
          calibrated effective density \num{2.5e17}\cmcu{} (not identifiable from $\Id$--$\Vg$) \lab{fitted}.
      \end{itemize}
    \end{column}
  \end{columns}
  \takeaway{Thin \InO{} switches off; confinement, surfaces and few donors then shape the device.}
\end{frame}

% ---------------------------------------------------------------------------------------------------------------------
\begin{frame}{W keeps \InO{} mobile; its role in nanometre-thin films is open}
  \oneTop
  {\scriptsize\renewcommand{\arraystretch}{1.0}%
  % plain tabular (not tabularx): tabularx typesets its body several times and would repeat the footnote citations.
  % Footnote [1] holds the three W:In2O3 film and transistor studies (rows 1-3), [2] Swallow, [3] the device paper.
  \begin{tabular}{@{}>{\raggedright\arraybackslash}p{2.95cm}>{\raggedright\arraybackslash}p{2.45cm}%
      >{\raggedright\arraybackslash}p{\dimexpr\textwidth-5.4cm-4\tabcolsep\relax}@{}}
    \toprule
    Work & System & Finding used here \\\midrule
    Newhouse et al.\ 2005\srefs{Newhouse2005,Kizu2014,Kim2024} & W:\InO{} (IWO) films & high Hall mobility at
      substantial carrier density \\
    Kizu et al.\ 2014\footnotemark[1] & a-In--W--O TFTs & a few \% W: high mobility and stability at low process
      temperature \\
    Kim et al.\ 2024\footnotemark[1] & IWO TFTs & stability depends on channel thickness \\
    Swallow et al.\ 2020\sref{Swallow2020} & Mo:\InO{} & Mo 4d states high in the conduction band (resonant doping) \\
    Januar et al.\ 2026\sref{Januar2026} & IWO, about 2\,\% W & 2.0--13.2\,nm devices; analytic model fitted per
      device; crystalline PBE cells of \InO{} and W:\InO{} \\
    \bottomrule
  \end{tabular}\par}
  \vspace{0.5mm}
  \begin{itemize}
    \item The device paper\footnotemark[3] (pp.~3, 4, 7) reports that W passivates oxygen vacancies, and relates the
      fall of the peak mobility from about 27.4 to 5.1\cmVs{} (13.2 to 2.0\,nm) to a roughly fourfold rise of the trap
      density.
    \item Apart from that paper's crystalline cells, the project's literature search found no peer-reviewed DFT study
      of W in \InO{}; those cells did not test the vacancy-passivation statement.
  \end{itemize}
  \takeaway{Whether W passivates vacancies in \InO{} is open; Parts~8 and 9 add crystalline tests.}
\end{frame}

% ---------------------------------------------------------------------------------------------------------------------
\begin{frame}{The 2.0--13.2\,nm threshold step ($-0.58$\,V) is fitted, not explained}
  \oneTop
  \begin{columns}[T,onlytextwidth]
    \begin{column}{0.32\textwidth}
      \vspace{1mm}
      {\small
      \begin{tabular}{@{}r r r@{}}
        \toprule
        $t$ (nm) & $\Vthcc$ (V) & $\Delta\Vthcc$ (V) \\\midrule
        2.0  & 0.662 &  \\
        6.3  & 0.644 & \num{-0.018} \\
        13.2 & 0.083 & {\color{ntumaroon}\num{-0.561}} \\\midrule
        \multicolumn{2}{@{}l}{2.0\,$\to$\,13.2\,nm} & \num{-0.579} \\
        \bottomrule
      \end{tabular}\par}
      \vspace{0.5mm}\lab{derived}
      \src{Constant-current threshold at $\Id = 10^{-9}$\,A/\textmu m, $\Vd = 0.7$\,V (registered extraction);
        $\Delta$: change from the next thinner film; last row: the 2.0--13.2\,nm step. One device per film. A 31.8\,nm
        curve, supplied with the data but not part of the published 2.0--13.2\,nm series, never falls below
        $10^{-9}$\,A/\textmu m. Data: Prof.\ M.-H.\ Lee's group, NTU (device paper, footnote~1).}
    \end{column}
    \begin{column}{0.66\textwidth}
      \begin{itemize}
        \item \textbf{The 2.0--13.2\,nm step} is $\Delta\Vthcc = -0.58$\,V. Of it, $-0.56$\,V lies between 6.3 and
          13.2\,nm (\textbf{the 6.3--13.2\,nm drop}); from 2.0 to 6.3\,nm the threshold changes by only $-0.018$\,V
          \lab{derived}.
        \item The device paper\sref{Januar2026} fits eight parameters per device. This project's per-film models match
          each threshold and on-current by construction \lab{fitted}, yet miss the 6.3\,nm on-state shape and the
          13.2\,nm swing \lab{computed}.
        \item All three out-of-sample predictions of the 6.3\,nm threshold failed, missing by $-0.29$, $-0.37$ and
          $-0.23$\,V \lab{falsified}.
        \item A shared model needs material inputs instead: chiefly $\dEc(t)$ (affinity, hence $\Vth$), $m^{*}(t)$
          (density of states, band mobility) and whether W is ionised (its share of the carriers).
      \end{itemize}
    \end{column}
  \end{columns}
  \takeaway{The question: which of these inputs can first principles supply to one shared model?}
\end{frame}

% ---------------------------------------------------------------------------------------------------------------------
\begin{frame}{This work: computed inputs, a shared model, registered tests}
  \oneTop
  {\centering
  \begin{tikzpicture}[x=1mm,y=1mm,font=\scriptsize,line width=0.5pt,
      box/.style={draw,align=center,inner sep=1.2mm,minimum height=8mm},
      arr/.style={-{Stealth[length=2mm]},line width=0.6pt}]
    \node[box,text width=44mm] (meas) at (23.2,20) {\textbf{Measured transfer curves}\\4 films, 2.0--31.8\,nm};
    \node[box,text width=40mm] (extr) at (75,20) {\textbf{Registered extraction}\\$\Vthcc$, swing, mobility, floor};
    \node[box,text width=33mm] (cmp) at (122,20) {\textbf{Comparison}\\against registered criteria};
    \node[box,text width=44mm,align=left] (dft) at (23.2,4) {\textbf{First principles} (PBE; HSE06 checks)\\
      crystalline films 0.46--1.96\,nm\\ W and O vacancies in crystalline cells\\
      amorphous films 0.7--2.0\,nm (sealed)};
    \node[box,text width=40mm] (inp) at (75,4) {\textbf{Device inputs}\\$\dEc(t)$, $\dEg(t)$, $\Delta m^{*}(t)$:
      laws fitted to computed points;\\W taken as fully ionised \lab{inferred}};
    \node[box,text width=33mm] (tcad) at (122,4) {\textbf{Shared-parameter model}\\campaign C10, drift--diffusion;\\
      one set for 2.0, 6.3, 13.2\,nm};
    \draw[arr] (meas) -- (extr);
    \draw[arr] (extr) -- (cmp);
    \draw[arr] (dft) -- (inp);
    \draw[arr,ntumaroon] (inp) -- (tcad);
    \draw[arr] (tcad) -- (cmp);
  \end{tikzpicture}\par}
  \src{Project workflow. Maroon arrow: quantities transferred from first principles to the device model.}
  {\footnotesize\textbf{Legend.} Device campaign C10 runs in stages S0 to S7 (Part~4); results from stage S2 on are
    sealed.\par}
  \begin{itemize}\footnotesize
    \item Five fitted parameters are shared by the 2.0, 6.3 and 13.2\,nm films (a sixth only in the percolation-mobility
      branch); no film-specific charge or mobility.
    \item Tests were registered before their outcomes where possible (Part~3): register revision~2 was frozen on
      10 October 2026, after stages S0 and S1 had been seen. No holdout device exists, so all comparisons are
      retrospective.
  \end{itemize}
  \takeaway{First principles supplies the confinement laws, not the donor density or traps.}
\end{frame}
